% Nouveau EGA V (cote 25): archive cote 25, site page 48 (PDF page 49).
% typescript. Draft compared: none.
% AI-transcribed by Claude Opus 5.5 (claude-opus-5-5, adaptive thinking, effort high), Message Batches API, batch ega25b (msgbatch_01A7n5YuXubUPNK1zsepi8RG), 2026-11-13; not yet checked by eye.
\scanpage{25/48}
\mspage{\unsure{45|46}}


grâce au fait que $\underline{X'}$ sing est de dim $r-1$, où $r=\dim \check{P}$ !

que $\phi_{\hand{\zeta}}$ soit $X_{k(\hand{\zeta})}$-régulière et $Y_{\hand{\zeta}}$ satisfasse la condition $R_k$ géométri-%

que. Alors:

a) Pour que $Z_{k-1}$ soit de codimension $\geq 2$ dans $\overset{\hand{\vee}}{P}$ , il  s que
$X'_i$ de $X'$ de dim $\leq k$ soit \uline{réglée} pour $f$. 
toute composante irréductible 

b) Pour que $Z_k$ soit de codimension $\geq 2$ dans $\check{P}$ , il  s que

toute composante irréductible 
$X_i$ de dim $\leq k+1$ de $X$ soit formée de pts lisses de $X$ et soit \uline{réglée} pour $f$.

\unsure{i.e. toute composante irréductible $X'_i$ de $X'$ de}

 \unsure{de dim $\leq k+1$ est,  réglée rel. à $f$,  telle que $X'_i = X_i$.}
\unsure{soit contenue} \unsure{dans $X'$}

En effet, pour tout $\hand{\zeta}$, l'ensemble géométriquement singulier

de $Y_{\hand{\zeta}}$  (NB on se borne aux $\hand{\zeta}$ tels que $\phi_{\hand{\zeta}}$ soit $X_{k(\hand{\zeta})}$-

régulier, ce qui est inoffensif vu 8.2.) est la réunion de

sing$(Y'_{\hand{\zeta}}$ ) et de l'image inverse $T_{\hand{\zeta}}$ de $T$ dans $Y_{\hand{\zeta}}$ , donc la codimension
l'inf de
\placedmargin{\unsure{lieu non-ordin. géométrique}}
de cet ensemble singulier dans $Y_{\hand{\zeta}}$ est égale à \hand{(}codim (sing$(Y'_{\hand{\zeta}})$,$Y'_{\hand{\zeta}}$ )

et de codim $(T_{\hand{\zeta}},Y_{\hand{\zeta}})$. Bornons nous aux $\hand{\zeta}$ tels que \hand{)}
\placedmargin{$*$} sing$(Y'_\zeta)$ soit fini
\hand{/}(ce qui est inoffensif, car  cela conduit à se mettre dans le

complémentaire d'un ensemble de codimension $\geq 2$\hand{\unsure{$\zeta$}$\}$}), Les points singulier

géométriques de $Y'_{\hand{\zeta}}$ sont alors isolés. La conclusion résulte facilement de

là , et de 8.4. 

Combinant 8.13.  8.5.et 8.6.\hand{/}, on trouve à la façon habituelle:

\uline{Corollaire} 8.14. On suppose $f:X\rightarrow P$ non ramifié et $X$ sans pts isolés.

\placedmargin{\unsure{n}} a) Supposons $X$ séparable sur $k$. Pour que l'ensemble  des $\hand{\zeta \in}$ $\overset{\hand{\vee}}{P}$

tels que $\phi_{\hand{\zeta}}$ soit $X_{k(\hand{\zeta})}$ -régulier et $Y_{\hand{\zeta}}$ soit séparable ait un complé-%

mentaire de codimension $\geq 2$, il f et s qu

\mspage{48}
